\documentclass[10pt,a4paper]{amsart}
\usepackage[margin=19mm]{geometry}
\usepackage{amsmath,amssymb,mathtools}
\usepackage[scr=rsfs,cal=euler]{mathalfa}
\usepackage{xfrac}
\usepackage[unicode,hidelinks,bookmarks=false]{hyperref}
\newcommand{\ie}{i.e.\ }
\newcommand{\cf}{cf.\ }
\newcommand{\resp}{respectively\ }
\newcommand{\Subsection}{\subsection}
\providecommand{\nobreakdash}{\nobreak\hbox{-}}
\setlength{\emergencystretch}{2em}
\multlinegap0pt
\usepackage{bm}
\usepackage{bbm}
\usepackage{dsfont}
\usepackage{xspace}
\usepackage{cancel}
\usepackage{mathtools}
\usepackage{amsthm}
\makeatletter
\@ifundefined{thm}{\newtheorem{thm}{Theorem}}{}
\makeatother
\newcommand{\bn}{\beta}
\newcommand{\dep}{depth}
\renewcommand{\geq}{\geqslant}
\renewcommand{\leq}{\leqslant}
\renewcommand{\le}{\leqslant}
\title[Shallow brambles]{Shallow brambles}
\author{Nicolas Bousquet, Wouter Cames van Batenburg, Louis Esperet, Gwena\"el Joret, Piotr Micek}
\date{Source: arXiv:2502.04177v4 (CC BY 4.0)}
\begin{document}
\maketitle

\noindent\emph{Source and licence.} Shallow brambles. Version arXiv:2502.04177v4 (CC BY 4.0), distributed under the Creative Commons Attribution 4.0 International licence, as recorded in the arXiv licence metadata for this exact version and shown on the abstract page https://arxiv.org/abs/2502.04177v4. Full rights notice: licenses/arxiv-2502.04177-CC-BY-4.0.txt.\par\smallskip

\noindent\textit{Editorial scope.} Counted unit: the Thm whose source label is \texttt{None} in the article (source0.tex lines 380--383, source label \texttt{None}), delivered together with the article's own complete original proof of it (source0.tex lines 384--422, one \texttt{proof} environment of 39 lines). No mathematics is written, repaired or re-derived: statement and proof are the article's own text line for line. Required context retained uncounted: none. Every definition, display and label of those lines is retained unchanged. Omitted from this package and not redistributed: 1--379, 423--617. Nothing inside a retained excerpt is omitted or altered: every retained body is delivered exactly as the article's lines are written, so no omission receipt of any kind accompanies this package. Citations: the retained excerpts carry no citation command at all, so no bibliography entry of the article is transcribed and no reference is redistributed. Reference adaptation: none was needed and none was made; every reference inside the delivered unit resolves inside it, so no unresolved reference or citation is produced. Printed slips kept literal: no printed word, symbol, hypothesis or number of the counted statement or of its proof was altered. Layout and class: the article's own class is \texttt{dmtcs-episciences} with packages including inputenc, subfigure, amsmath, amssymb, amsthm, amscd, verbatim, cleveref, mathtools, hyperref, enumitem; the wrapper is the lane's pinned \texttt{amsart} editorial wrapper at 10pt on A4, which restates the theorem-like environments the retained text needs and adds the article's own macro definitions that the retained text needs (\texttt{\textbackslash bn}, \texttt{\textbackslash dep}, \texttt{\textbackslash geq}, \texttt{\textbackslash le}, \texttt{\textbackslash leq}), with extra wrapper packages (bm, bbm, dsfont, xspace, cancel, mathtools). The delivered numbering is the unit's own. Assets: no retained span contains a figure environment, a graphics-inclusion command, a PDF-page inclusion, a table or a TikZ picture; no image file is copied into this package, the package declares no asset dependency and no image file is redistributed with it. Licence: version arXiv:2502.04177v4 of this article is distributed under the Creative Commons Attribution 4.0 International licence, as recorded in the arXiv licence metadata for this exact version and shown on the abstract page https://arxiv.org/abs/2502.04177v4. A full rights notice is delivered at licenses/arxiv-2502.04177-CC-BY-4.0.txt. Characters: every retained excerpt is pure ASCII, so the delivered engine, XeLaTeX with its default Latin Modern text font, reports no missing character.
\begin{thm}
For every graph $G$ and integer $r\geq0$, 
 $$  \bn_r(G) \leq (5r+1)\cdot (\omega_{5r+1}(G))^2.$$
\end{thm}

\begin{proof}
 We prove that $\bn_r(G) \leq (5r+1) \cdot t^2$, where $t:=\omega_{5r+1}(G)$.
Let $\mathcal{B}$ be a \dep-$r$ bramble of $G$.
We call $\mathcal{M}$ a \emph{good model} if
\begin{itemize}
\item $\mathcal{M}$ is a depth-$(5r+1)$ minor-model of a complete graph in $G$, and
\item $|V(M)|\leq 1+5(r+1)(t-1)$ for every branch set $M\in \mathcal{M}$.
\end{itemize}

Note that there exists at least one good model, since a single vertex defines one. Given a minor-model $\mathcal{M}$, let $V(\mathcal{M})$ denote the union of the branch sets of $\mathcal{M}$, and let $R(\mathcal{M})$ denote the union of all bramble elements of $\mathcal{B}$ that are disjoint from $V(\mathcal{M})$. Observe that, since every two bramble elements touch, the subgraph induced by $R(\mathcal{M})$ has radius at most $3r+1$.

\smallskip

From now on, we fix a good model $\mathcal {M}$ in $G$ such that
$(|R(\mathcal{M})|, |V(\mathcal{M})|)$ is lexicographically minimized among all good models. Let $M_1,M_2,\ldots, M_s$ be the branch sets of $\mathcal{M}$. Note that by the definition of $t$, we have $s\leq t$.

\smallskip

The proof is done if $R(\mathcal{M})= \emptyset$, because then $V(\mathcal{M})$ is a hitting set of $\mathcal{B}$ of order at most \[\sum_{1=1}^{s} |V(M_i)| \leq t+ (5r+1)(t-1)t \leq (5r+1)\cdot t^2.\]
Thus we may assume that $R(\mathcal{M})$ is non-empty, and from there we will derive a contradiction.
\smallskip


For every branch set $M_i\in \mathcal{M}$ it holds that
 \begin{equation}\label{ineq:}
 |R(\mathcal{M}-M_i)| > |R(\mathcal{M})|, 
 \end{equation}
  for otherwise we would have $|R(\mathcal{M}-M_i)| \le|R(\mathcal{M})|$ and $|V(\mathcal{M}-M_i)| <|V(\mathcal{M})|$, contradicting the minimality of $\mathcal{M}$. 



Thanks to (\ref{ineq:}) we know that each branch set $M_i\in \mathcal{M}$ satisfies the following key
property: There is a bramble element $B\in\mathcal{B}$ that intersects $M_i$ and avoids all other branch sets of $\mathcal{M}$. Take a shortest path $P_i$ from $M_i$ to the non-empty set $R(\mathcal{M})$ through $M_i \cup B$. Let $x_i$ be the endpoint of $P_i$ in $R(\mathcal{M})$. Observe that $P_i$ contains at most $2r+1$ edges, and thus at most $2r+2$ vertices. Let $Q_i$ be $P_i$ minus its endpoint in $M_i$ (so $Q_i$ contains at most $2r$ edges).

Recall that the subgraph induced by $R(\mathcal{M})$ has radius at most $3r+1$, so there exists a vertex $x^*\in R(\mathcal{M})$ and for each $1\le i \le s$, a path on at most $3r+1$ edges, included in $R(\mathcal{M})$ and connecting $x^*$ to  $x_i$. Let $X\subseteq R(\mathcal{M})$ be the union of the vertex sets of these $s$ paths. Note that $X$ is connected and contains at most $1+(3r+1)s$ vertices. 

Now define a new branch set $M_{s+1}$, consisting of $X$ together with the paths $Q_1,\ldots, Q_s$. Then the radius of $M_{s+1}$ is at most $5r+1$, and $|V(M_{s+1})|\leq 1+(5r+1)s\le 1+(5r+1)(t-1)$. 
Moreover, by construction $M_{s+1}$ is connected to all previous branch sets. We conclude that $\mathcal{M}+M_{s+1}$ is a good model with $|R(\mathcal{M}+M_{s+1})|<|R(\mathcal{M})|$, contradicting the minimality of $\mathcal{M}$.
\end{proof}

\end{document}
